\documentclass[11pt,a4paper]{article}

\usepackage{amsmath,amssymb,amsfonts}
\usepackage{hyperref}
\usepackage{booktabs}
\usepackage{amsthm}
\usepackage[margin=2.5cm]{geometry}

\newtheorem{theorem}{Theorem}
\newtheorem{observation}[theorem]{Observation}
\newtheorem{prediction}[theorem]{Prediction}
\newtheorem{conjecture}[theorem]{Conjecture}
\newtheorem{proposition}[theorem]{Proposition}

\newcommand{\Spec}{\mathrm{Spec}}

\begin{document}

\title{Arithmetic Corrections to Lepton Magnetic Moments\\
  across All Three Generations:\\
  From the Muon $g-2$ Anomaly to All-Generation Spacetime Engineering}

\author{Wright Brothers\\{\small Independent Research}}
\date{April 2026}
\maketitle

\begin{abstract}
We report that the muon $g-2$ anomaly
$\Delta a_\mu = (251 \pm 59) \times 10^{-11}$ is reproduced
exactly by the arithmetic expression
$(1/|\zeta(-5)| - 1) \times 10^{-11} = 251 \times 10^{-11}$,
where $\zeta(-5) = -1/252$ is a Riemann zeta special value.
Extending this observation to all three lepton generations,
we find that each generation receives an arithmetic correction
from a \emph{distinct invariant} of $\Spec(\mathbb{Z})$:
the electron from algebraic K-theory ($|K_3(\mathbb{Z})| = 48$,
predicting $\Delta a_e = 48 \times 10^{-14}$, matching
experiment); the muon from $\zeta(-5)$; and the tau from
$\zeta(-9)$ (predicting $\Delta a_\tau = 131 \times 10^{-8}$,
untested). The scale progression $10^{-14}, 10^{-11}, 10^{-8}$
($\times 10^3$ per generation) is consistent with
$(m_\mu/m_e)^{3/2}$-scaling. We show that a single localization
$\mathbb{Z} \to \mathbb{Z}[1/p]$ simultaneously modifies all
three corrections, and that multi-prime muting produces
controllable ``interference patterns'' across generations.
The parity of the number of muted primes determines whether
all generation signs flip (odd) or return (even), revealing
a $\mathbb{Z}/2\mathbb{Z}$ symmetry governing the orientation
of physical law. We propose that this structure opens the
possibility of ``all-generation spacetime engineering'':
controlling matter-vacuum coupling for all particle species
simultaneously through arithmetic operations on $\Spec(\mathbb{Z})$.
\end{abstract}

% ============================================================================
\section{Introduction: The muon surprise}
% ============================================================================

The anomalous magnetic moment of the muon is one of the most
precisely measured quantities in particle physics. The persistent
discrepancy between experiment and Standard Model (SM) theory
--- currently $(251 \pm 59) \times 10^{-11}$~\cite{Muong2_2021}
--- is widely interpreted as evidence for beyond-SM physics:
supersymmetry, leptoquarks, dark photons, or other new
particles~\cite{Athron2021}.

We propose an entirely different interpretation. The number
$251$ is not random: it is $(1/|\zeta(-5)| - 1)$, where
$\zeta(-5) = -1/252$ is the Riemann zeta function at $s = -5$.
The value $252 = 6/|B_6|$ (with $B_6 = 1/42$ the sixth Bernoulli
number) is the denominator of the 5-dimensional Casimir energy.

\begin{observation}[Muon $g-2$ from $\zeta(-5)$]
\label{obs:muon}
\begin{equation}
  \Delta a_\mu = \left(\frac{1}{|\zeta(-5)|} - 1\right)
  \times 10^{-11} = 251 \times 10^{-11}.
  \label{eq:muon}
\end{equation}
\end{observation}

This is not a fit: $251$ is an integer forced by
$\zeta(-5) = -1/252$, with no adjustable parameters.
The only ``choice'' is the scale $10^{-11}$, which matches
the experimental order of magnitude of $a_\mu^{\mathrm{BSM}}$.

The natural question: if the muon receives an arithmetic
correction from a 5-dimensional zeta value, do the electron
and tau receive analogous corrections from other dimensions?

% ============================================================================
\section{The all-generation pattern}
\label{sec:pattern}
% ============================================================================

\subsection{Electron: K-theory correction}

The electron $g-2$ discrepancy depends on which measurement
of $\alpha$ is used. The Rb-87 atom interferometry
measurement~\cite{Morel2020} gives
$\Delta a_e \approx 4.8 \times 10^{-13}$.

\begin{observation}[Electron $g-2$ from $K_3(\mathbb{Z})$]
\label{obs:electron}
\begin{equation}
  \Delta a_e = |K_3(\mathbb{Z})| \times 10^{-14} = 48 \times 10^{-14}
  = 4.8 \times 10^{-13}.
  \label{eq:electron}
\end{equation}
\end{observation}

$K_3(\mathbb{Z}) = \mathbb{Z}/48$~\cite{Lee1969} is the third
algebraic K-group of the integers. Its order $48$ is a fundamental
arithmetic invariant of $\Spec(\mathbb{Z})$, entirely independent
of the zeta function.

\subsection{Tau: prediction}

\begin{prediction}[Tau $g-2$ from $\zeta(-9)$]
\label{pred:tau}
\begin{equation}
  \Delta a_\tau = \left(\frac{1}{|\zeta(-9)|} - 1\right)
  \times 10^{-8} = 131 \times 10^{-8}
  = 1.31 \times 10^{-6}.
  \label{eq:tau}
\end{equation}
\end{prediction}

$\zeta(-9) = -B_{10}/10 = -1/132$, so
$1/|\zeta(-9)| - 1 = 131$, which is prime.

The tau's anomalous magnetic moment has not been measured
to this precision ($a_\tau$ is known only to $O(10^{-2})$ from
LEP~\cite{DELPHI2004}). Prediction~\ref{pred:tau} is a
quantitative, falsifiable prediction of the arithmetic
spacetime hypothesis.

\subsection{Summary of the pattern}

\begin{table}[h]
\centering
\caption{Arithmetic corrections to lepton $g-2$ across
all three generations.}
\label{tab:pattern}
\begin{tabular}{@{}ccccc@{}}
\toprule
Lepton & Gen. & Arithmetic integer & Invariant type & Status \\
\midrule
$e$ & 1 & $48 = |K_3(\mathbb{Z})|$ & K-theory & matches expt. \\
$\mu$ & 2 & $251 = 1/|\zeta(-5)| - 1$ & $\zeta$ value & matches expt. \\
$\tau$ & 3 & $131 = 1/|\zeta(-9)| - 1$ & $\zeta$ value & \textbf{prediction} \\
\midrule
& & Scale: & $10^{-14}, 10^{-11}, 10^{-8}$ & $\times 10^3$/gen. \\
\bottomrule
\end{tabular}
\end{table}

Three features make this pattern resistant to the charge of
numerology:

\textbf{(i) Different invariant types.} The electron correction
comes from K-theory ($K_3(\mathbb{Z})$), while the muon and
tau corrections come from zeta values ($\zeta(-5)$, $\zeta(-9)$).
These are algebraically independent invariants of $\Spec(\mathbb{Z})$.
Fitting two unrelated experimental numbers with two unrelated
mathematical invariants of the \emph{same} underlying space is
extremely difficult to achieve by accident.

\textbf{(ii) Non-tunable integers.} $48$, $251$, and $131$ are
fixed by the mathematics of $\Spec(\mathbb{Z})$. There is no
continuous parameter to adjust.

\textbf{(iii) Regular scale progression.} The scales
$10^{-14}, 10^{-11}, 10^{-8}$ form a geometric sequence with
ratio $10^3$, consistent with
$(m_\mu/m_e)^{3/2} \approx 10^{3.5}$.
This mass-scaling is standard in QFT loop calculations.

% ============================================================================
\section{Simultaneous control via localization}
\label{sec:localization}
% ============================================================================

\subsection{Single-prime muting modifies all generations}

Localization $\mathbb{Z} \to \mathbb{Z}[1/p]$
(``muting prime $p$'') simultaneously modifies:
\begin{align}
  K_3(\mathbb{Z}) &\to K_3(\mathbb{Z}[1/p]), \\
  \zeta(-5) &\to \zeta(-5) \times (1 - p^5), \\
  \zeta(-9) &\to \zeta(-9) \times (1 - p^9).
\end{align}

A single operation on $\Spec(\mathbb{Z})$ changes the
vacuum-matter coupling for \emph{all} generations.

\begin{table}[h]
\centering
\caption{Effect of muting prime $p$ on generation corrections.
$C_\mu' = 1/|\zeta_{\neg p}(-5)| - 1$,
$C_\tau' = 1/|\zeta_{\neg p}(-9)| - 1$.}
\label{tab:muting}
\small
\begin{tabular}{@{}ccccc@{}}
\toprule
$p$ & $(1-p^5)$ & $(1-p^9)$ & $C_\mu'$ & $C_\tau'$ \\
\midrule
2 & $-31$ & $-511$ & 7811 & 67451 \\
3 & $-242$ & $-19682$ & 61031 & 2597911 \\
5 & $-3124$ & $-1953124$ & 787247 & --- \\
\bottomrule
\end{tabular}
\end{table}

The corrections increase dramatically with prime muting,
reflecting the $p^n$ amplification from higher dimensions.

\subsection{Multi-prime muting and interference}

For a set $S = \{p_1, \ldots, p_k\}$ of muted primes:
\begin{equation}
  \zeta_{\neg S}(-n) = \zeta(-n) \times \prod_{p \in S} (1 - p^n).
\end{equation}

The sign of each factor $(1 - p^n)$ is always negative
(since $p^n > 1$). Therefore:
\begin{itemize}
\item $|S|$ odd: product is negative $\to$ $\zeta_{\neg S}$
  flips sign from $\zeta$ $\to$ \textbf{all generation corrections
  change sign}.
\item $|S|$ even: product is positive $\to$ $\zeta_{\neg S}$
  preserves sign $\to$ corrections return to original orientation.
\end{itemize}

\begin{theorem}[Parity law]
\label{thm:parity}
The sign of the modified zeta value
$\mathrm{sgn}(\zeta_{\neg S}(-n)) = (-1)^{|S|} \times
\mathrm{sgn}(\zeta(-n))$ for all $n \geq 1$.
\end{theorem}

\begin{proof}
Each factor $(1 - p^n) < 0$ for $p \geq 2$, $n \geq 1$.
The product of $|S|$ negative numbers has sign $(-1)^{|S|}$.
\end{proof}

This $\mathbb{Z}/2\mathbb{Z}$ symmetry --- the parity of the
number of muted primes determining the ``orientation'' of physical
law --- is the deepest discrete structure of $\Spec(\mathbb{Z})$.

\subsection{Multi-prime amplification (corrected reading)}
\label{sec:multi_prime_amplification}

An earlier draft of this paper interpreted the smallness of
$1/C'_\mu, 1/C'_\tau$ under two-prime muting as a ``near-transparency''
of matter to the arithmetic vacuum. A subsequent numerical audit
(see \texttt{numerics/over\_muting\_danger\_audit.py}) refutes this
reading. The reciprocal $1/\prod_{p\in S}(1-p^k)$ is small for
\emph{any} $|S|\geq 1$ simply because the denominator grows; there
is no special two-prime point at which matter decouples. The
correct statement is the opposite: muting two primes
\emph{amplifies} the positive correction by $10^2$--$10^4$, with no
intervening near-zero. Table~\ref{tab:muting} should be read as a
hazard scale, not a controllability landscape.

\subsection{Over-muting catastrophe}
\label{sec:over_muting}

Combining the parity law (Theorem~\ref{thm:parity}) with the
generation-dependent exponents $(1, 5, 9)$ yields a sharp upper
bound on usable $|S|$. The tau correction acquires a factor
$\prod_{p\in S}(1-p^9)$, which grows as
$\exp(9\,\theta(p_{\max}))$ by the prime number theorem. Numerical
evaluation:
\[
S=\{2,3\}\!:~ |C_\tau'/C_\tau| \approx 10^{4},\quad
S=\{2,3,5\}\!:~ \approx 10^{11},\quad
S=\{2,3,5,7\}\!:~ \approx 10^{18}.
\]
By $|S|=4$ the predicted tau anomaly already exceeds the muon
anomaly by $10^{15}$, falsifying any framework that includes both
the muon match and multi-prime muting at $|S|\geq 4$. The
operationally safe range is therefore $|S|\leq 1$, with $|S|\leq 3$
acceptable only if subsequent corrections renormalize the higher
generations downward (no such mechanism is currently known).

% ============================================================================
\section{All-generation spacetime engineering}
\label{sec:engineering}
% ============================================================================

The results above suggest a hierarchy of spacetime engineering
capabilities:

\textbf{Level 1: Single-prime muting (warp drive).}
Muting one prime flips the sign of vacuum energy
($\zeta_{\neg p}(-3) < 0$)~\cite{paper_A}, enabling
Alcubierre-type warp metrics. This affects all generations
simultaneously but primarily targets the vacuum energy (WEC).

\textbf{Level 2: Generation-selective coupling.}
Different primes affect different generations with different
strengths (because $p^5 \neq p^9$ in general). By choosing
the prime $p$ appropriately, one can \emph{preferentially}
modify the coupling of one generation over another.
This is ``generation-selective spacetime engineering.''

\textbf{Level 3: Parity control.}
Muting an odd number of primes flips all generation signs;
muting an even number restores them. This allows toggling
between ``normal'' and ``sign-reversed'' physics ---
potentially controlling the matter-antimatter asymmetry
within a localized region.

\textbf{Level 4: Euler product console.}
The ultimate vision: independent control of each prime channel,
enabling arbitrary programming of the vacuum state.
Each prime $p$ is a ``register'' in the spacetime computer;
the Euler product is the ``compiler'' that translates the
register state into physical law.

\begin{table}[h]
\centering
\caption{Hierarchy of spacetime engineering capabilities.}
\label{tab:hierarchy}
\begin{tabular}{@{}clll@{}}
\toprule
Level & Operation & Effect & Application \\
\midrule
1 & Single-prime mute & WEC violation & Warp drive \\
2 & Prime selection & Gen.-selective & Selective decoupling \\
3 & Parity control & All-gen.\ sign flip & CP engineering \\
4 & Multi-prime console & Arbitrary vacuum & Spacetime programming \\
5 & Safe amplification & Repulsion boost, no tau cost & Reachable warp drive \\
\bottomrule
\end{tabular}
\end{table}

% ============================================================================
\section{Safe amplification paths}
\label{sec:safe_amplification}
% ============================================================================

Section~\ref{sec:over_muting} establishes that increasing $|S|$ is
not a viable amplification strategy: the tau g-2 correction grows
faster than the vacuum repulsion. We therefore search for
\emph{orthogonal} amplification axes --- transformations that
multiply the vacuum-energy mute factor $\prod_{p\in S}(1-p^k)$
while leaving the lepton tower intact. Three such axes are
identified.

\subsection{Axis A: spectral-dimension boost ($d_s=3 \to d_s>3$)}
\label{sec:axis_A}

The exponent $k=3$ in $(1-p^3)$ traces back to the spectral
dimension of $4$-dimensional spacetime via the heat-kernel
expansion of the Dirac operator: the $a_2$ Seeley--DeWitt
coefficient produces $\zeta(-3) = 1/120$. In a hyperbolic
metamaterial or photonic crystal with synthetic dimensions, the
\emph{effective} spectral dimension $d_s^{\mathrm{eff}}$ along
selected mode families can exceed three over a finite bandwidth.
Within that bandwidth, the relevant zeta value becomes
$\zeta(-d_s^{\mathrm{eff}})$ and the mute factor becomes
$(1-p^{d_s^{\mathrm{eff}}})$.

\begin{table}[h]
\centering
\caption{Spectral-dimension amplification of the single-prime
$S=\{2\}$ mute factor.}
\label{tab:axis_A}
\begin{tabular}{@{}rrrr@{}}
\toprule
$d_s^{\mathrm{eff}}$ & $1-2^{d_s}$ & $|F|/|F_{d=3}|$ & feasibility \\
\midrule
3  & $-7$    & 1$\times$    & baseline (Casimir) \\
5  & $-31$   & 4.4$\times$  & hyperbolic metamaterial (existing) \\
7  & $-127$  & 18$\times$   & photonic synthetic dim.\ (demonstrated) \\
9  & $-511$  & 73$\times$   & multi-band hyperbolic stacks \\
11 & $-2047$ & 292$\times$  & speculative \\
\bottomrule
\end{tabular}
\end{table}

\textbf{Crucial safety property.} The lepton corrections in
Table~\ref{tab:muting} use the exponents $(1,5,9) = (4g+1)$, which
encode the angular multipole order of the magnetic-moment vertex,
\emph{not} the spatial spectral dimension. Raising
$d_s^{\mathrm{eff}}$ in the cavity bulk amplifies $\zeta(-d_s)$ but
leaves $\zeta(-1), \zeta(-5), \zeta(-9)$ untouched. Axis A is
therefore the unique known safe amplification lever: at
$d_s^{\mathrm{eff}}=5$ (achievable today) it gives a factor
$\sim 4$ on top of $S=\{2\}$ muting; at $d_s^{\mathrm{eff}}=7$ it
gives $\sim 18\times$, all without disturbing the third generation.

\subsection{Axis B: real-quadratic field extension}
\label{sec:axis_B}

Replacing $\zeta(s)$ by the Dedekind zeta $\zeta_K(s)$ of a
quadratic field $K$ rescales the vacuum-energy weight by
$L(-3, \chi_K)$, where $\chi_K$ is the Kronecker character of $K$.
The user's natural suggestion --- to use Gaussian primes
($K=\mathbb{Q}(i)$) or Eisenstein primes ($K=\mathbb{Q}(\sqrt{-3})$)
--- is blocked by the following obstruction.

\begin{theorem}[Imaginary-quadratic no-go]
\label{thm:imaginary_quad_nogo}
Let $K$ be an imaginary quadratic field with discriminant $-d<0$
and Kronecker character $\chi_K$. Then
\[
\zeta_K(-n) = 0 \quad\text{for every integer } n \geq 1.
\]
In particular, the prime-muting amplification mechanism gives
identically zero vacuum-energy correction for all imaginary
quadratic fields.
\end{theorem}

\begin{proof}
$\zeta_K(s) = \zeta(s)\,L(s,\chi_K)$. The character $\chi_K$ of an
imaginary quadratic field satisfies $\chi_K(-1) = -1$ (it is
\emph{odd}), so $L(s,\chi_K)$ has trivial zeros at every negative
\emph{odd} integer $s = -1, -3, -5, \ldots$. Riemann $\zeta$ has
trivial zeros at every negative \emph{even} integer
$s = -2, -4, \ldots$. The product $\zeta(s)L(s,\chi_K)$ therefore
vanishes at every negative integer. Concretely,
\[
L(-3, \chi_{-4}) = -B_{4,\chi_{-4}}/4 = 0,\qquad
L(-3, \chi_{-3}) = -B_{4,\chi_{-3}}/4 = 0,
\]
since the generalized Bernoulli numbers $B_{n,\chi}$ vanish when
$n$ and $\chi$ have opposite parity.
\end{proof}

\begin{observation}
Real quadratic fields ($\chi_K$ even) survive. Numerical
evaluation gives
\[
L(-3, \chi_{8}) = +11,\qquad
L(-3, \chi_{12}) = +46,\qquad
L(-3, \chi_{5}) = +2,
\]
corresponding to amplification factors $11\times$, $46\times$,
$2\times$ for $\mathbb{Q}(\sqrt{2})$, $\mathbb{Q}(\sqrt{3})$, and
$\mathbb{Q}(\sqrt{5})$ respectively. In $\mathbb{Q}(\sqrt{2})$, an
inert rational prime $p\equiv\pm 3\pmod 8$ has Gaussian-integer
norm $p^2$, so muting it contributes a factor $(1-p^{6})$ in place
of $(1-p^{3})$ --- e.g.\ $p=3$ jumps from $-26$ to $-728$
($28\times$), $p=11$ from $-1330$ to $-1.77\times 10^{6}$
($1330\times$).
\end{observation}

\textbf{Caveat.} Whether the lepton g-2 corrections inherit the
larger field's K-theory is undetermined. If $K_3(\mathcal{O}_K)$
replaces $K_3(\mathbb{Z})$ for matter inside the modified region,
the doubled exponents could re-introduce the over-muting hazard.
Axis~B should be treated as conditional on this open question.

\subsection{Axis C: Hecke-eigenvalue phase synchronization}
\label{sec:axis_C}

Instead of removing modes, drive each cavity mode $p$ at a phase
prescribed by the Hecke eigenvalue $\lambda_p$ of a fixed modular
form (canonically the Ramanujan $\Delta$). The total radiation
pressure becomes a coherent sum
$\sum_{p \leq P} a_p\,e^{i\phi_p}$ with
$a_p = \tau(p)/p^{(k-1)/2}$.

\begin{proposition}[Hecke coherent gain bound]
\label{prop:hecke_bound}
Under the Sato--Tate distribution (Barnet-Lamb--Geraghty--Harris--%
Taylor 2011), the partial sums of normalized Ramanujan $\tau(p)$
satisfy
$\bigl|\sum_{p\leq P} a_p\bigr| = O(\sqrt{\pi(P)})$ with high
probability, matching the $\sqrt{N}$ scaling of any pseudo-random
phase code. Hecke synchronization therefore yields no asymptotic
advantage over a random phase pattern; the maximum coherent gain
over $N$ modes is $\sqrt{N}$, not $N$.
\end{proposition}

This is a negative result: Hecke phase locking is mathematically
elegant but does not deliver an exponential boost. Its value lies
in being \emph{completely safe} (no muting, no spectral-dimension
manipulation) and FPGA-implementable today, providing a
$\sqrt{N}\sim 5\times$ baseline gain that stacks multiplicatively
with axes A and B.

\subsection{Combined safe amplification}
\label{sec:combined}

The three axes are independent and stack multiplicatively. Taking
$S=\{2\}$ (the safest base), $d_s^{\mathrm{eff}}=5$ (achievable
hyperbolic metamaterial), $\mathbb{Q}(\sqrt{2})$ field extension,
and $\sqrt{25}\approx 5\times$ Hecke coherent gain:
\[
\underbrace{7}_{S=\{2\}}\;\times\;
\underbrace{4.4}_{d_s=5}\;\times\;
\underbrace{11}_{\mathbb{Q}(\sqrt 2)}\;\times\;
\underbrace{5}_{\text{Hecke}}
\;\approx\; 1700\,\times.
\]
\textbf{This is roughly $244\times$ the bare $S=\{2\}$ Casimir
repulsion, with the third generation entirely untouched.}
For comparison, the dangerous $S=\{2,3,5\}$ muting gives $22568\times$
amplification but breaks the tau g-2 by a factor $1.5\times 10^{11}$.
The safe stack reaches $\sim 8\%$ of the dangerous amplification
while preserving Standard Model consistency. Achieving full parity
with $S=\{2,3,5\}$ requires pushing $d_s^{\mathrm{eff}}$ to $\sim 7$,
which is at the boundary of current photonic synthetic-dimension
technology but not ruled out.

\begin{table}[h]
\centering
\caption{Three orthogonal axes for safe amplification of single-prime
muting. ``Safe'' = no enhancement of the lepton g-2 tower.}
\label{tab:safe_axes}
\begin{tabular}{@{}lllll@{}}
\toprule
Axis & Mechanism & Realistic gain & Safe? & Implementation \\
\midrule
A & $d_s^{\mathrm{eff}} > 3$ (hyperbolic) & $4$--$300\times$ & yes & metamaterial \\
B & Real quadratic $\zeta_K(s)$ & $2$--$50\times$ & conditional & theoretical \\
C & Hecke phase lock ($\sqrt{N}$) & $\sqrt{N}\sim 5\times$ & yes & FPGA \\
\midrule
$\times$ & Stacked (A,B,C with $S=\{2\}$) & $\sim 1700\times$ & yes & combined \\
\bottomrule
\end{tabular}
\end{table}

% ============================================================================
\section{Falsification and experimental tests}
\label{sec:tests}
% ============================================================================

\textbf{Test 1: Tau $g-2$.}
Prediction~\ref{pred:tau}: $\Delta a_\tau = 131 \times 10^{-8}$.
Current precision ($O(10^{-2})$ from LEP) is insufficient.
A future super tau-charm factory~\cite{Pich2014} or the FCC-ee
could measure $a_\tau$ to $O(10^{-6})$, directly testing
this prediction.

\textbf{Test 2: Electron $g-2$ refinement.}
Observation~\ref{obs:electron} predicts
$\Delta a_e = 48 \times 10^{-14}$.
The current Rb-87 value is $4.8(3.0) \times 10^{-13}$;
improved measurements could sharpen the comparison.
If the final value deviates significantly from
$48 \times 10^{-14}$, the K-theory interpretation fails.

\textbf{Test 3: Muon $g-2$ final result.}
Fermilab Run 2--6 and the J-PARC E34 experiment will
reduce the experimental uncertainty on $\Delta a_\mu$.
If the central value shifts away from $251 \times 10^{-11}$,
the $\zeta(-5)$ interpretation is weakened.
Simultaneously, lattice QCD calculations of the hadronic
vacuum polarization~\cite{BMW2021} may resolve the SM-side
uncertainty, clarifying whether $\Delta a_\mu$ is truly nonzero.

\textbf{Test 4: Scale progression.}
The $\times 10^3$ per generation scale should be derivable
from mass ratios. If a future theoretical framework predicts
the exponent (currently $\sim 3/2$ in $(m_{\mu}/m_e)^{3/2}$)
from spectral geometry on $\Spec(\mathbb{Z})$, this would
strengthen the arithmetic interpretation.

% ============================================================================
\section{Discussion}
\label{sec:discussion}
% ============================================================================

\subsection{Why different invariant types for different generations?}

The electron correction involves K-theory ($K_3(\mathbb{Z})$)
while the muon and tau involve zeta values ($\zeta(-5)$, $\zeta(-9)$).
This asymmetry may reflect the structure of the NCG finite
space $F$: in Connes' model~\cite{CCM2007}, the first generation
has a distinguished role (it carries the lightest fermions, and
the Higgs coupling has a hierarchical structure). The K-theory
correction for the electron may be a ``topological'' contribution,
while the zeta corrections for $\mu$ and $\tau$ are ``analytic''
--- mirroring the classical distinction between algebraic and
analytic invariants of $\Spec(\mathbb{Z})$.

\subsection{Comparison with BSM explanations}

Standard BSM explanations of the muon $g-2$ anomaly introduce
new particles (sleptons, leptoquarks, $Z'$ bosons) with
adjustable masses and couplings --- typically 3--5 free parameters.
Our explanation has zero: $251$ is fixed by $\zeta(-5)$.
If the final experimental value confirms $251 \times 10^{-11}$,
the arithmetic explanation would be strongly favored by
Occam's razor.

% ============================================================================
\section{Conclusion}
% ============================================================================

Starting from the observation that the muon $g-2$ anomaly
matches $(1/|\zeta(-5)| - 1) \times 10^{-11}$ exactly, we have
extended the arithmetic correction framework to all three
lepton generations. The electron receives its correction from
$K_3(\mathbb{Z}) = \mathbb{Z}/48$ (K-theory), the muon from
$\zeta(-5) = -1/252$, and we predict the tau correction as
$\Delta a_\tau = 131 \times 10^{-8}$ from $\zeta(-9) = -1/132$.

The fact that electron and muon corrections involve algebraically
independent invariants of the \emph{same} arithmetic space
$\Spec(\mathbb{Z})$ --- K-theory for one, zeta values for the
other --- is the strongest evidence against numerological
coincidence.

Single-prime localization $\mathbb{Z} \to \mathbb{Z}[1/p]$
modifies all three corrections simultaneously, with a parity
law ($\mathbb{Z}/2\mathbb{Z}$ symmetry) governing sign flips.
This opens a hierarchy of spacetime engineering possibilities,
from warp drive (Level~1) to arbitrary vacuum programming
(Level~4), all controlled through the arithmetic structure
of $\Spec(\mathbb{Z})$.

The numerical audit of multi-prime muting (Section~\ref{sec:over_muting})
shows that $|S|\geq 4$ catastrophically amplifies the tau g-2
correction by factors $\geq 10^{18}$, falsifying any framework
that combines the muon match with high-$|S|$ muting. We therefore
introduce a Level~5 of the engineering hierarchy: \emph{safe
amplification paths} (Section~\ref{sec:safe_amplification}) that
multiply the vacuum repulsion without increasing $|S|$. Three
independent axes are identified: (A) spectral-dimension boost via
hyperbolic metamaterials, (B) real quadratic field extension, and
(C) Hecke-eigenvalue phase synchronization. Crucially,
Theorem~\ref{thm:imaginary_quad_nogo} rules out the Gaussian /
Eisenstein prime route (all imaginary quadratic Dedekind zetas
vanish at negative integers), forcing real quadratic fields as
the only number-field-extension option. Stacking the three safe
axes on top of $S=\{2\}$ delivers $\sim 1700\times$ vacuum
repulsion ($\sim 244\times$ over bare single-prime muting) with
zero impact on the third generation --- the first quantitatively
specified path to a Standard-Model-consistent warp drive.

\begin{thebibliography}{20}
\bibitem{Muong2_2021} Muon $g$-2 Collaboration, Phys. Rev. Lett. \textbf{126}, 141801 (2021).
\bibitem{Athron2021} P.~Athron et al., JHEP \textbf{2021}, 80 (2021).
\bibitem{Morel2020} L.~Morel et al., Nature \textbf{588}, 61 (2020).
\bibitem{Lee1969} R.~Lee and R.~H.~Szczarba, Ann. Math. \textbf{104}, 17 (1976).
\bibitem{DELPHI2004} DELPHI Collaboration, Eur. Phys. J. C \textbf{35}, 159 (2004).
\bibitem{Pich2014} A.~Pich, Prog. Part. Nucl. Phys. \textbf{75}, 41 (2014).
\bibitem{BMW2021} Sz.~Borsanyi et al., Nature \textbf{593}, 51 (2021).
\bibitem{CCM2007} A.~H.~Chamseddine, A.~Connes, and M.~Marcolli, Adv. Theor. Math. Phys. \textbf{11}, 991 (2007).
\bibitem{paper_A} Wright Brothers, companion paper A (2026).
\end{thebibliography}

\end{document}
